% Standalone source. Compile with pdflatex or latexmk -pdf.
\documentclass[11pt]{article}
\usepackage[T1]{fontenc}
\usepackage{lmodern}
\usepackage[margin=1in]{geometry}
\usepackage{amsmath,amssymb,amsthm,mathtools}
\usepackage{microtype,enumitem,booktabs,longtable,array,xcolor,listings}
\usepackage{tikz}
\usetikzlibrary{arrows.meta,positioning}
\usepackage{xurl}
\usepackage[colorlinks=true,linkcolor=blue!45!black,urlcolor=blue!45!black]{hyperref}
\usepackage{fancyhdr}
\pagestyle{fancy}\fancyhf{}
\fancyhead[L]{\small De Rham instance: Lean blueprint}
\fancyhead[R]{\small\thepage}
\setlength{\headheight}{14pt}
\setlength{\emergencystretch}{2em}
\setlist{itemsep=3pt,topsep=5pt}
\lstset{basicstyle=\ttfamily\footnotesize,breaklines=true,columns=fullflexible,keepspaces=true,frame=single,rulecolor=\color{black!20},xleftmargin=4pt,xrightmargin=4pt}
\newcommand{\A}{\mathbb A}
\newcommand{\PP}{\mathbb P}
\newcommand{\Gm}{{\mathbb G_m}}
\newcommand{\Z}{\mathbb Z}
\newcommand{\N}{\mathbb N}
\newcommand{\cA}{\mathcal A}
\newcommand{\cP}{\mathcal P}
\newcommand{\OS}{\operatorname{OS}}
\newcommand{\pOS}{\overline{\operatorname{OS}}}
\newcommand{\im}{\operatorname{im}}
\newcommand{\kerop}{\operatorname{ker}}
\newcommand{\res}{\operatorname{res}}
\newcommand{\codim}{\operatorname{codim}}
\newcommand{\Hom}{\operatorname{Hom}}
\newcommand{\Spec}{\operatorname{Spec}}
\newcommand{\SH}{\operatorname{SH}}
\newcommand{\one}{\mathbf 1}
\newcommand{\logc}[1]{[#1]}
\newcommand{\planned}[1]{\texttt{\detokenize{#1}}}
\newtheorem{theorem}{Theorem}[section]
\newtheorem{lemma}[theorem]{Lemma}
\newtheorem{proposition}[theorem]{Proposition}
\newtheorem{corollary}[theorem]{Corollary}
\theoremstyle{definition}
\newtheorem{definition}[theorem]{Definition}
\theoremstyle{remark}
\newtheorem{remark}[theorem]{Remark}

\newcommand{\lean}[1]{\par\smallskip{\footnotesize\color{blue!50!black}\textsf{Planned declaration: }\nolinkurl{#1}}\par\smallskip}
\newcommand{\uses}[1]{}

\newcommand{\Om}{\Omega}
\newcommand{\dR}{\mathrm{dR}}
\newcommand{\dd}{\mathrm d}
\newcommand{\Der}{\operatorname{Der}}
\newcommand{\gr}{\operatorname{gr}}
\newcommand{\LL}{L}
\newcommand{\extlabel}[1]{\texttt{\detokenize{#1}} in the first blueprint}
\hypersetup{pdftitle={Algebraic de Rham cohomology as an ordinary theory: Lean blueprint},pdfauthor={Prepared for Adam Topaz and the formalization team}}
\title{A de Rham instance of the cohomology axioms\\\large A second Lean formalization blueprint}
\author{Prepared for Adam Topaz and the formalization team}
\date{6 November 2026}
\begin{document}
\maketitle
\begin{abstract}
We construct the requested ordinary cohomology theory on every smooth algebra over a characteristic-zero field. The proof includes a basis-free descent construction of the de Rham differential, polynomial homotopy invariance, the smooth-divisor residue sequence, and a global contracting homotopy on every positive pole-order quotient. No coordinate charts, arrangements, or flatness of morphisms of pairs enter the construction. We verify the exact signs and naturality required by the fixed Lean interface, including units, disconnected algebras, and the zero ring. The existing arrangement structure theorem is then instantiated, not reproved.

All declarations in the namespace \texttt{OrlikSolomon.DeRham} are proposed. Nothing in this document has been formally verified or elaborated in Lean. The fixed interface in the team's prompt is taken as authoritative; no field of it is weakened.
\end{abstract}
\paragraph{Interface verdict.}
All fields listed in the prompt hold for this model. One example in the prompt needs a factual correction: \(k[z]\to k[x,y,y^{-1}]\), \(z\mapsto xy\), is flat (indeed a Laurent-polynomial extension after changing coordinates). We keep it as the unit-change test and supply a genuinely non-flat variant in the regression section.
\clearpage
\tableofcontents
\newpage
\section{Scope, conventions, and the construction target}

Write $A_f=A[1/f]$ and $S=A/(f)$. Initially the base $k$ is an arbitrary commutative ring; the point at which $k$ must contain $\mathbb Q$ will be explicit. The final instance takes $k$ to be a field of characteristic zero. All algebra maps are unital $k$-algebra maps, including maps to the zero ring. All de Rham differentials and complexes are $k$-linear, not $A$-linear. Extend complexes by zero in negative degrees.

The coefficient variable $t$ is a formal Laurent variable, independent of any geometric coordinate. Set
\[
 B=k[t,t^{-1}],\qquad \mathcal B_w=k\,t^w,\qquad
 C(A)=\left(\bigoplus_{p\geq0}H^p_{\dR}(A/k)\right)\otimes_k B.
\]
Its component of bidegree $(p,w)$ is $H^p_{\dR}(A/k)t^w$. Coefficients have cohomological degree zero, so the ordinary algebra tensor product here has the correct multiplication. Denote the untwisted cohomological residue by $r$ and the untwisted Gysin map by $G$; the interface maps are
\[
 \operatorname{residue}(x t^w)=r(x)t^{w-1},\qquad
 \operatorname{gysin}(y t^w)=G(y)t^{w+1}.
\]

For a cochain complex $K$, our shift convention is $(K[-1])^n=K^{n-1}$ with differential $-d_K$. Residues will put $df/f$ on the \emph{left}. Thus the unshifted form residue anticommutes with $d$, while the residue into $\Om_S^\bullet[-1]$ is a chain map. Cohomology of this shifted complex is identified with ordinary cohomology of $\Om_S$ by the identity on cocycles and boundaries, with no degree-dependent scalar.

The main construction is first made for the canonical quotient and localization. The prompt permits the implementers to transport it to a \planned{DivisorPair} whose closed and complement algebras are merely isomorphic to these. The transport maps are determined by its specified $i$ and $j$; they are not arbitrary isomorphisms. All naturality statements below are expressed in terms of these maps and therefore survive transport.

\section{Exterior algebra facts used in the construction}

\begin{lemma}[Exterior powers, base change, and quotients]\label{lem:dr-exterior}
\lean{OrlikSolomon.DeRham.exteriorBaseChangeAndQuotient}
Let $M$ be an $A$-module, $N\subseteq M$ a submodule, and $A\to A'$ a commutative algebra map. There are natural identifications
\[
 A'\otimes_A\bigwedge_A^nM\simeq\bigwedge_{A'}^n(A'\otimes_AM),\qquad
 \bigwedge_A^n(M/N)\simeq
 \frac{\bigwedge_A^nM}{N\wedge\bigwedge_A^{n-1}M}.
\]
The denominator denotes the submodule spanned by the indicated wedges. Exterior words span, each degree-one element squares to zero, and homogeneous elements of degrees $p,q$ commute with sign $(-1)^{pq}$, over any base ring.
\end{lemma}
\begin{proof}
The base-change map sends $a'\otimes(m_1\wedge\cdots\wedge m_n)$ to $a'(1\otimes m_1)\wedge\cdots\wedge(1\otimes m_n)$. Its inverse sends $(a_1'\otimes m_1)\wedge\cdots\wedge(a_n'\otimes m_n)$ to $(\prod a_i')\otimes(m_1\wedge\cdots\wedge m_n)$. Balancedness and alternation verify both maps; they are inverse on spanning words. For the quotient, the wedge of the classes of $m_i$ is unchanged modulo the indicated submodule if one lift is changed by an element of $N$. This gives the inverse to the natural surjection and proves the claimed kernel. Spanning follows from the tensor-algebra presentation. Expanding $(u+v)^2=0$ gives $uv=-vu$, and moving letters proves the sign rule. These arguments use the exterior relation $u^2=0$ itself and never divide by two.
\end{proof}

\begin{lemma}[Contraction by a functional]\label{lem:dr-contraction}
\lean{OrlikSolomon.DeRham.exteriorContraction}
For an $A$-linear map $\lambda:M\to A$, there is a degree $-1$ operator $\iota_\lambda$ on $\bigwedge_A M$ with
\[
 \iota_\lambda(m_1\wedge\cdots\wedge m_n)
 =\sum_{j=1}^n(-1)^{j-1}\lambda(m_j)
 m_1\wedge\cdots\widehat{m_j}\cdots\wedge m_n.
\]
It is $A$-linear and satisfies
\[
 \iota_\lambda(\alpha\wedge\beta)
 =\iota_\lambda\alpha\wedge\beta+(-1)^p\alpha\wedge\iota_\lambda\beta,
 \qquad
 \iota_\lambda(e\wedge\alpha)+e\wedge\iota_\lambda\alpha=\lambda(e)\alpha.
\]
If $\lambda(e)=1$, then
\[
 \ker(e\wedge-:\textstyle\bigwedge^n M\to\bigwedge^{n+1}M)
 =e\wedge\textstyle\bigwedge^{n-1}M.
\]
\uses{lem:dr-exterior}
\end{lemma}
\begin{proof}
Define the signed deletion operator on the tensor algebra. It is compatible with tensor balancing because $\lambda$ is $A$-linear. On a defining exterior relation $m\otimes m$ it gives $\lambda(m)m-m\lambda(m)=0$; on a product containing such a relation, the signed product rule places the result in the exterior ideal. Thus it descends, giving the displayed formula and product rule. Apply that rule to $e\wedge\alpha$ for the second identity. If $e\wedge\alpha=0$ and $\lambda(e)=1$, it gives $\alpha=e\wedge\iota_\lambda\alpha$. The reverse inclusion follows from $e\wedge e=0$. The zero and negative degrees are interpreted using zero modules.
\end{proof}

\section{A de Rham dg algebra without a basis of differentials}

The operator $d$ is not $A$-linear. Defining it by an $A$-linear exterior-algebra lift of $a\mapsto da$ would be incorrect. We instead construct it on a canonical polynomial presentation and descend through a differential ideal. This uses a basis only for polynomial differentials, never for $\Om^1_{A/k}$.

\begin{definition}[Forms and universal differential]\label{def:dr-forms}
\lean{OrlikSolomon.DeRham.forms}
For any commutative $k$-algebra set $\Om_A^0=A$, $\Om_A^n=\bigwedge_A^n\Om^1_{A/k}$ for $n>0$, and $\Om_A^n=0$ for $n<0$. Here $d_0:A\to\Om^1_{A/k}$ is the universal $k$-derivation. The direct sum of these modules is the exterior algebra, with wedge multiplication. Every form is a finite sum of expressions $a_0\,da_1\wedge\cdots\wedge da_n$.
\end{definition}
\begin{proof}
The module of K\"ahler differentials is generated by the symbols $da$ with additivity, $d(c)=0$ for $c\in k$, and $d(ab)=a\,db+b\,da$. Applying the exterior spanning assertion gives the last claim. Equivalently derive it from the universal presentation of the K\"ahler module, so it applies to mathlib's chosen implementation.
\end{proof}

\begin{lemma}[Polynomial de Rham differential]\label{lem:dr-polynomial-dga}
\lean{OrlikSolomon.DeRham.polynomialDGA}
For $P=k[X_i\mid i\in I]$, even with $I$ infinite, $\Om_P^1$ is free with basis $dX_i$. There is a $k$-linear operator of degree one determined by
\[
 d(p\,dX_{i_1}\wedge\cdots\wedge dX_{i_n})
 =\sum_{j\in I}\partial_j(p)\,dX_j\wedge dX_{i_1}\wedge\cdots\wedge dX_{i_n}.
\]
The sum has finite support. It satisfies $d^2=0$ and the signed Leibniz rule.
\uses{def:dr-forms,lem:dr-exterior}
\end{lemma}
\begin{proof}
A derivation $P\to M$ is uniquely determined by its values on the $X_i$: on a monomial the Leibniz rule gives the usual finite sum, and extending $k$-linearly constructs it for arbitrary assigned values. This proves the claimed free description by the universal property of $\Om_P^1$.

For a precise definition of $d$, identify
\[
 \Om_P^n\simeq P\otimes_k\bigwedge_k^n(k^{(I)}).
\]
The linear map $p\mapsto\sum_j\partial_j(p)\otimes e_j$ has finite support and is $k$-linear. Wedge its second factor on the left with an exterior word; the resulting bilinear operation descends through the tensor product. This defines the displayed operator on every element, without choosing an order on $I$. The product rule for partial derivatives, and moving $dX_j$ past $n$ letters, give
$d(\alpha\beta)=d\alpha\,\beta+(-1)^n\alpha\,d\beta$.

Finally $d^2$ on a spanning word has coefficients $\partial_l\partial_jp$ in front of $dX_l\wedge dX_j\wedge dX_I$. The terms $l=j$ vanish and the terms $(l,j)$ and $(j,l)$ cancel. The partial derivatives commute because on each monomial both orders multiply by the same two integer exponents (or by $a(a-1)$ when equal). This also proves the assertion in positive characteristic. Only finitely many variables occur in each calculation.
\end{proof}

\begin{lemma}[Polynomial presentation and the differential ideal]\label{lem:dr-presentation}
\lean{OrlikSolomon.DeRham.presentationDifferentialIdeal}
For an arbitrary $A$, let $P_A=k[X_a\mid a\in A]$, let $\epsilon:P_A\twoheadrightarrow A$ send $X_a$ to $a$, and put $I=\ker\epsilon$. The homogeneous ideal
\[
 J^n=I\Om_{P_A}^n+
 \operatorname{span}_{P_A}\{dh\wedge\eta:h\in I,\ \eta\in\Om_{P_A}^{n-1}\}
\]
is stable under $d$. Its graded algebra quotient is canonically $\bigoplus_n\Om_A^n$.
\uses{lem:dr-polynomial-dga,lem:dr-exterior}
\end{lemma}
\begin{proof}
This is the homogeneous two-sided ideal generated by $I$ in degree zero and $dI$ in degree one: graded commutativity moves a degree-one generator to the front of a word. For $h\in I$,
\[
 d(h\omega)=dh\wedge\omega+h\,d\omega,\qquad
 d(dh\wedge\eta)=-dh\wedge d\eta,
\]
so it is a differential ideal.

The module $(A\otimes_{P_A}\Om_{P_A}^1)/(1\otimes dh:h\in I)$ represents $k$-derivations from $A$. Indeed a derivation from $P_A$ into an $A$-module factors through $A$ exactly when it kills $I$, and these are exactly the indicated relations. Its universal property identifies it with $\Om_A^1$. Exterior base change and the exterior quotient formula of Lemma~\ref{lem:dr-exterior} now identify the degree-$n$ quotient with $\Om_A^n$. In particular the degree-zero quotient is $A$. No smoothness, flatness, or basis of $\Om_A^1$ has been used.
\end{proof}

\begin{definition}[De Rham complex and dg algebra]\label{def:dr-complex}
\lean{OrlikSolomon.DeRham.complex}
Transport the differential descended in Lemma~\ref{lem:dr-presentation} to $\Om_A^\bullet$. Denote it by $d$. With the wedge product this is the de Rham dg algebra of $A/k$.
\uses{lem:dr-presentation}
\end{definition}

\begin{theorem}[The differential and functoriality]\label{thm:dr-dga}
\lean{OrlikSolomon.DeRham.dgaLawsAndFunctoriality}
For every commutative $k$-algebra,
\[
 d(a_0\,da_1\wedge\cdots\wedge da_n)
 =da_0\wedge da_1\wedge\cdots\wedge da_n.
\]
The operator squares to zero and satisfies the signed Leibniz rule. Wedge multiplication is associative, unital, and graded commutative. Every $k$-algebra map $\phi:A\to A'$ induces a dg algebra map $\phi_*$, sending $a_0\,da_1\cdots da_n$ to $\phi(a_0)d\phi(a_1)\cdots d\phi(a_n)$; these maps respect identity and composition.
\uses{def:dr-complex,def:dr-forms,lem:dr-exterior}
\end{theorem}
\begin{proof}
Lift the $a_i$ to the canonical variables. The formula follows from the polynomial dg algebra and its quotient. The two differential identities descend through the differential ideal, and all multiplicative identities descend from the exterior algebra. For functoriality, the K\"ahler universal property gives the semilinear map $da\mapsto d\phi(a)$, and exterior functoriality gives its maps in every degree. On the spanning expressions in the theorem, applying $d$ before or after this map gives the same expression. Thus it is a chain map. Identity and composition are checked on those same expressions. The formula also shows uniqueness of the differential with these properties, so the transported operator has no residual dependence on a presentation or on auxiliary exterior identifications.
\end{proof}

\begin{lemma}[Localization of forms and its differential]\label{lem:dr-localization}
\lean{OrlikSolomon.DeRham.localizationForms}
For any $A$ and $f\in A$, there are natural identifications
\[
 \Om_{A_f}^n\simeq (\Om_A^n)_f,
 \qquad d(f^{-m}\omega)=f^{-m}d\omega-m f^{-m-1}df\wedge\omega
 \quad(m\geq0).
\]
In particular $d(df/f)=0$.
\uses{thm:dr-dga,lem:dr-exterior}
\end{lemma}
\begin{proof}
A derivation $\delta:A\to N$ into an $A_f$-module has a unique extension, necessarily
\[
 \delta(a/f^m)=f^{-m}\delta(a)-m a f^{-m-1}\delta(f).
\]
This is well defined: if $a/f^m=b/f^n$, some $f^r(af^n-bf^m)=0$ in $A$. Differentiate this equality into $N$, invert $f$, and use $af^n-bf^m=0$ in $A_f$ to obtain equality of the two displayed values. The sum and product rules follow by a common-denominator calculation. The universal property of differentials therefore gives $\Om_{A_f}^1\simeq A_f\otimes_A\Om_A^1$, and exterior base change gives all degrees. The differential formula follows from the signed Leibniz rule and $d(f^{-1})=-f^{-2}df$. Finally $d(df/f)=-f^{-2}df\wedge df=0$ by the exterior relation. These assertions do not require that localization be injective.
\end{proof}

\begin{lemma}[Forms on a principal quotient]\label{lem:dr-quotient}
\lean{OrlikSolomon.DeRham.principalQuotientForms}
For $S=A/(f)$ the restriction maps are surjective in each degree, with
\[
 \Om_S^n\simeq
 \Om_A^n/(f\Om_A^n+df\wedge\Om_A^{n-1}).
\]
They commute with $d$. In particular there is a right exact conormal sequence
\[
 (f)/(f^2)\longrightarrow\Om_A^1\otimes_AS
 \longrightarrow\Om_S^1\longrightarrow0,
 \quad f\bmod f^2\longmapsto df\bmod f.
\]
\uses{thm:dr-dga,lem:dr-exterior}
\end{lemma}
\begin{proof}
A derivation $A\to N$ into an $S$-module factors through $S$ iff it kills $f$. Indeed $d(af)=a\,df+f\,da$, whose second term acts as zero in $N$. Thus the differential module of $S$ is $(\Om_A^1/f\Om_A^1)/(S\,df)$. Exterior base change and the quotient formula give the stated result in every degree. The map from $(f)/(f^2)$ is well defined because $d(f^2a)$ is zero modulo $f$. Its image is exactly $S\,df$. Functoriality of the de Rham complex gives the chain-map assertion. No left injectivity of the conormal map is claimed at this stage.
\end{proof}

\begin{definition}[Cohomology algebra]\label{def:dr-cohomology}
\lean{OrlikSolomon.DeRham.cohomologyAlgebra}
Define $H^n_{\dR}(A/k)=\ker(d:\Om_A^n\to\Om_A^{n+1})/\im(d:\Om_A^{n-1}\to\Om_A^n)$. Wedge induces a unital graded-commutative $k$-algebra structure on its direct sum, functorial covariantly in $A$.
\uses{thm:dr-dga}
\end{definition}
\begin{proof}
Products of cocycles are cocycles. If $\beta$ is closed, $(d\eta)\beta=d(\eta\beta)$; if $\alpha$ is closed of degree $p$, $\alpha\,d\eta=(-1)^p d(\alpha\eta)$. Thus changing either representative changes the product by a boundary. This gives the multiplication on the quotient; its ring identities and graded commutativity follow from forms. The class of $1$ is the unit. Chain maps preserve cycles, boundaries, and products, and hence induce the functor. Boundaries form an ideal in the algebra of cycles, not in the algebra of all forms; use that distinction in the Lean quotient construction.
\end{proof}

\section{Polynomial homotopy invariance}

\begin{lemma}[Polynomial decomposition of forms]\label{lem:dr-polynomial-splitting}
\lean{OrlikSolomon.DeRham.polynomialFormDecomposition}
For every $A$, writing the new variable as $z$, every form on $A[z]$ decomposes uniquely as
\[
 \omega=\alpha(z)+dz\wedge\beta(z),\qquad
 \alpha\in\Om_A^n[z],\quad\beta\in\Om_A^{n-1}[z].
\]
Here polynomial notation for a module means finite sums with coefficients in that module. The differential is
\[
 d\omega=d_A\alpha+dz\wedge(\partial_z\alpha-d_A\beta).
\]
\uses{thm:dr-dga,lem:dr-exterior}
\end{lemma}
\begin{proof}
A derivation from $A[z]$ is equivalent to a derivation from $A$ together with the value assigned to $z$, by the polynomial Leibniz formula. Hence
$\Om_{A[z]}^1\simeq(A[z]\otimes_A\Om_A^1)\oplus A[z]dz$.
In exterior degree $n$, expand a word using this direct sum, move $dz$ to the left, and discard terms with two copies of $dz$. The inverse sends the two summands to regular words and words with $dz$ first. The projections onto zero and one $dz$ factors show uniqueness. Applying $d$ to $z^m\eta$ and to $dz\wedge z^m\eta$ gives the displayed formula, with the minus sign from the degree of $dz$.
\end{proof}

\begin{theorem}[Integration homotopy]\label{thm:dr-homotopy}
\lean{OrlikSolomon.DeRham.polynomialHomotopyInvariance}
Suppose every positive integer is invertible in $k$. For every commutative $k$-algebra, the inclusion $A\to A[z]$ is a cohomology isomorphism, with inverse induced by evaluation at $z=0$. A chain homotopy from inclusion followed by evaluation to the identity is
\[
 K(\alpha+dz\wedge\beta)=I\beta,
 \qquad I\Bigl(\sum_{m\geq0}z^m\beta_m\Bigr)
 =\sum_{m\geq0}\frac{z^{m+1}}{m+1}\beta_m.
\]
\uses{lem:dr-polynomial-splitting,def:dr-cohomology}
\end{theorem}
\begin{proof}
The sum is finite, and the denominator is a scalar from $k$. Direct coefficient calculations give $\partial_zI=1$, $I\partial_z\alpha=\alpha-\alpha(0)$, and $d_AI=Id_A$. Therefore
\[
 \begin{aligned}
 dK(\omega)&=d_AI\beta+dz\wedge\beta,\\
 Kd(\omega)&=I\partial_z\alpha-Id_A\beta,
 \end{aligned}
 \qquad (dK+Kd)(\omega)=\omega-\alpha(0).
\]
Evaluation kills $dz$ and restricts $\alpha$ to $\alpha(0)$, so this is precisely $dK+Kd=1-\operatorname{inc}\operatorname{ev}_0$. The other composite is the identity already on forms. On cocycles the homotopy changes a representative by a boundary, proving that the induced maps are inverse. They are algebra homomorphisms by functoriality, so they give an algebra equivalence on total cohomology. This is the first use of rational scalars; smoothness has not been used.
\end{proof}

\section{The algebra of a smooth principal pair}

In this section $f$ is a non-zero-divisor in $A$, and both $A$ and $S=A/(f)$ are formally smooth over $k$. Finite presentation, part of the given \planned{SmoothAlg} predicate, is not needed by these arguments. No connectedness or nonzero-ring assumption is made.

\begin{lemma}[Projectivity and injectivity into localization]\label{lem:dr-projective}
\lean{OrlikSolomon.DeRham.formsProjectiveAndLocalizationInjective}
If $A$ is formally smooth over $k$, then $\Om_A^1$ and every $\Om_A^n$ are projective $A$-modules. If $f$ is a non-zero-divisor, multiplication by $f$ is injective on each $\Om_A^n$, and $\Om_A^n\to\Om_{A_f}^n$ is injective.
\uses{def:dr-forms,lem:dr-exterior,lem:dr-localization}
\end{lemma}
\begin{proof}
Let $N\twoheadrightarrow N'$ be a surjection of $A$-modules, and let $\delta:A\to N'$ be a $k$-derivation. The map $a\mapsto(a,\delta(a))$ is a $k$-algebra map $A\to A\ltimes N'$. The surjection $A\ltimes N\to A\ltimes N'$ has square-zero kernel, so formal smoothness lifts this map to $A\ltimes N$. Its first component is necessarily the identity, and its second component is a derivation lifting $\delta$. The universal property of $\Om_A^1$ therefore says that $\Hom_A(\Om_A^1,-)$ lifts through every surjection, which is projectivity.

A projective module is a retract of a free module. Apply $\bigwedge^n$ to a retraction to make $\bigwedge^n\Om_A^1$ a retract of the corresponding exterior power of a free module. The latter is free, with basis the increasing words in a basis of the original free module: ordering the basis and repeatedly using exterior anticommutation gives spanning, and coefficient alternating maps give independence. This includes infinite bases because each word is finite. Thus all exterior powers are projective; degree zero is $A$ itself.

Multiplication by $f$ is injective on a free module coordinatewise, and hence on a retract. Its iterates are also injective. An element of a module maps to zero in its localization at $f$ exactly when some power of $f$ kills it. This proves the final assertion using Lemma~\ref{lem:dr-localization}. The same reasoning applies to the zero module over the zero ring.
\end{proof}

\begin{lemma}[Split conormal sequence]\label{lem:dr-conormal}
\lean{OrlikSolomon.DeRham.splitConormalSequence}
Put $\overline M=\Om_A^1\otimes_AS$ and $e=df\bmod f\in\overline M$. There is an $S$-linear map $\lambda:\overline M\to S$ with $\lambda(e)=1$. In particular
\[
 0\longrightarrow S\xrightarrow{1\mapsto e}\overline M
 \longrightarrow\Om_S^1\longrightarrow0
\]
is split exact, and $\overline M\simeq S e\oplus\Om_S^1$ after a splitting is chosen.
\uses{lem:dr-quotient}
\end{lemma}
\begin{proof}
Multiplication by $f$ identifies $S$ with $(f)/(f^2)$: if $af=f^2b$, cancellation of $f$ gives $a=fb$. Consider the square-zero extension $A/(f^2)\twoheadrightarrow S$. Formal smoothness of $S$ lifts its identity to a $k$-algebra section $\sigma:S\to A/(f^2)$.

For $a\in A$ put $\delta(a)=a\bmod f^2-\sigma(a\bmod f)$, an element of $(f)/(f^2)$. The section is additive and respects $k$. Multiplying the two expressions $\sigma(\bar a)+\delta(a)$ and $\sigma(\bar b)+\delta(b)$, the product of the two kernel terms is zero, giving
\[
 \delta(ab)=\bar a\,\delta(b)+\bar b\,\delta(a).
\]
After identifying the kernel with $S$, this is a derivation $A\to S$ taking $f$ to $1$. Its universal linear map factors through $\overline M$ and is the required $\lambda$.

Right exactness and the description of the kernel come from Lemma~\ref{lem:dr-quotient}. The equality $\lambda(e)=1$ makes the first map injective. Explicitly $m\mapsto m-\lambda(m)e$ projects onto $\ker\lambda$ and identifies the quotient $\overline M/Se$ with $\ker\lambda$. This provides the full splitting without assuming a basis or constant rank. Here the non-zero-divisor hypothesis identifies the conormal line; smoothness of $S$ supplies the section. Smoothness of $A$ was not needed for this particular lemma.
\end{proof}

\begin{lemma}[A global normal derivation modulo the divisor]\label{lem:dr-normal-derivation}
\lean{OrlikSolomon.DeRham.existsNormalDerivationModulo}
There exist a $k$-derivation $v:A\to A$ and $h\in A$ such that
\[
 v(f)=1+fh.
\]
If $g=v(f)$, then $(f,g)=A$, and on $A_g$ the differential $df$ is unimodular: the extended derivation $v/g$ takes $f$ to $1$.
\uses{lem:dr-projective,lem:dr-conormal,lem:dr-localization}
\end{lemma}
\begin{proof}
Compose $\Om_A^1\to\overline M$ with $\lambda$ to obtain an $A$-linear map to $S$. Projectivity of $\Om_A^1$ lifts it through $A\twoheadrightarrow S$ to an $A$-linear map $\ell:\Om_A^1\to A$. Set $v=\ell\circ d_0$. Its value on $f$ is $1$ modulo $(f)$, so it equals $1+fh$ for some $h$. Thus $1=g-fh$, proving the unit-ideal assertion. Extend $v$ to $A_g$ by the localization formula and multiply it by $g^{-1}$ to obtain a derivation taking $f$ to $1$.

This gives a two-open cover $D(f)\cup D(g)=\Spec A$: on $D(f)$ the divisor is absent, while $df$ is unimodular on $D(g)$. We will not glue de Rham cohomology over this cover. All subsequent proofs use the single global derivation $v$ and the congruence $v(f)\equiv1\pmod f$. This is how unimodularity only along the divisor is sufficient, even when $df$ is not unimodular on all of $\Spec A$.
\end{proof}

\begin{lemma}[Wedge exactness modulo the divisor]\label{lem:dr-modf-koszul}
\lean{OrlikSolomon.DeRham.wedgeDfKernelModulo}
If $\eta\in\Om_A^q$ satisfies $df\wedge\eta\in f\Om_A^{q+1}$, then its image in $\Om_S^q$ is zero. Equivalently, modulo $f$, the kernel of $e\wedge-$ on $\bigwedge_S^q\overline M$ is the image of $e\wedge-$ in the preceding degree.
\uses{lem:dr-conormal,lem:dr-contraction,lem:dr-quotient,lem:dr-exterior}
\end{lemma}
\begin{proof}
Exterior base change identifies $\Om_A^q/f\Om_A^q$ with $\bigwedge_S^q\overline M$. The hypothesis becomes $e\wedge\bar\eta=0$. Contracting with the functional $\lambda$ from Lemma~\ref{lem:dr-conormal} gives $\bar\eta=e\wedge\iota_\lambda\bar\eta$. This becomes zero after quotienting by $e$, which is precisely the map to $\Om_S^q$. For $q=0$, the same contraction identity forces $\bar\eta=0$ directly. This is a global calculation in an $S$-module.
\end{proof}

\section{Logarithmic forms and the residue sequence}

From now on identify regular forms with their images in $\Om_{A_f}^\bullet$ using Lemma~\ref{lem:dr-projective}, and put $\theta=df/f$.

\begin{definition}[Logarithmic complex]\label{def:dr-log}
\lean{OrlikSolomon.DeRham.logComplex}
Define the $A$-submodule
\[
 \LL^n=\Om_A^n(\log f)=\Om_A^n+\theta\wedge\Om_A^{n-1}
 \subseteq\Om_{A_f}^n.
\]
In degree zero it is $A$, and it is zero in negative degrees.
\uses{lem:dr-projective,lem:dr-localization}
\end{definition}

\begin{lemma}[Log dg algebra]\label{lem:dr-log-dga}
\lean{OrlikSolomon.DeRham.logSubcomplexAndProduct}
The modules $\LL^n$ form a dg subalgebra of $\Om_{A_f}^\bullet$. In particular they are a graded right and left module over regular forms, and
\[
 d(\omega+\theta\wedge\eta)=d\omega-\theta\wedge d\eta.
\]
The graded ideal $f\LL^\bullet$ is also a subcomplex, so $\LL^\bullet/f\LL^\bullet$ has an induced differential.
\uses{def:dr-log,thm:dr-dga,lem:dr-localization}
\end{lemma}
\begin{proof}
The displayed formula follows from $d\theta=0$. Expanding the product of two log forms, the term with two $\theta$ factors vanishes after moving them together, and each remaining term is regular or has exactly one $\theta$ factor. Thus products stay in $\LL$. Finally $df=f\theta$ and
\[
 d(f\xi)=df\wedge\xi+f\,d\xi=f(\theta\wedge\xi+d\xi)\in f\LL.
\]
This verifies differential stability of the quotient even though $d$ is not $A$-linear.
\end{proof}

\begin{lemma}[Regularity criterion for logarithmic terms]\label{lem:dr-regularity}
\lean{OrlikSolomon.DeRham.logTermRegularIffRestrictionZero}
For $\eta\in\Om_A^{n-1}$,
\[
 \theta\wedge\eta\in\Om_A^n
 \quad\Longleftrightarrow\quad \eta|_S=0.
\]
\uses{lem:dr-modf-koszul,lem:dr-quotient,lem:dr-projective}
\end{lemma}
\begin{proof}
If $\theta\wedge\eta=\omega$ is regular, multiply by $f$ in the localization. Injectivity on regular $n$-forms makes this the equality $df\wedge\eta=f\omega$ in $\Om_A^n$. Lemma~\ref{lem:dr-modf-koszul} gives $\eta|_S=0$. Conversely the quotient description gives
$\eta=f\eta_1+df\wedge\eta_2$, allowing finite sums to be combined into this form by linearity. Then
$\theta\wedge\eta=df\wedge\eta_1$ is regular. The $n=1$ case simply has $\eta_2=0$, and the nonpositive degrees follow from the zero conventions.
\end{proof}

\begin{definition}[Residue of forms]\label{def:dr-residue}
\lean{OrlikSolomon.DeRham.formResidue}
Define $\rho:\LL^n\to\Om_S^{n-1}$ by
\[
 \rho(\omega+\theta\wedge\eta)=\eta|_S.
\]
It is $A$-linear with $A$ acting on the target through $S$, and hence $k$-linear.
\uses{lem:dr-regularity}
\end{definition}
\begin{proof}
If two decompositions represent the same log form, their difference has $\theta\wedge(\eta-\eta')$ regular. Lemma~\ref{lem:dr-regularity} says that their restrictions agree. Additivity and scalar compatibility follow by adding or multiplying decompositions, so the definition is independent of every choice.
\end{proof}

\begin{theorem}[Short exact residue sequence, with its sign]\label{thm:dr-residue-sequence}
\lean{OrlikSolomon.DeRham.logResidueShortExact}
There is a degreewise short exact sequence of complexes of $k$-modules
\[
 0\longrightarrow\Om_A^\bullet\longrightarrow\LL^\bullet
 \xrightarrow{\rho}\Om_S^\bullet[-1]\longrightarrow0.
\]
Equivalently, on the unshifted target, $\rho d=-d\rho$.
\uses{def:dr-residue,lem:dr-log-dga,lem:dr-quotient,lem:dr-projective,lem:dr-regularity}
\end{theorem}
\begin{proof}
Injectivity on the left is already established. Every $\gamma\in\Om_S^{n-1}$ lifts to a regular $\eta$, and $\theta\wedge\eta$ has residue $\gamma$, proving surjectivity. If a log form has zero residue, its logarithmic part is regular by Lemma~\ref{lem:dr-regularity}; hence the kernel is exactly the regular forms. Finally
$\rho d(\omega+\theta\wedge\eta)=-d(\eta|_S)=-d\rho(\omega+\theta\wedge\eta)$.
This is exactly the chain-map condition for the differential $-d$ on the shifted target. In degree zero the target is zero, as required.
\end{proof}

\begin{lemma}[Right projection formula on forms]\label{lem:dr-form-projection}
\lean{OrlikSolomon.DeRham.formResidueRightProduct}
For $\alpha\in\LL^n$ and $\beta\in\Om_A^q$,
\[
 \rho(\alpha\wedge\beta)=\rho(\alpha)\wedge(\beta|_S).
\]
There is no sign in this formula.
\uses{def:dr-residue,lem:dr-log-dga}
\end{lemma}
\begin{proof}
Write $\alpha=\omega+\theta\wedge\eta$. Its product with $\beta$ is
$\omega\wedge\beta+\theta\wedge(\eta\wedge\beta)$, so its residue is $(\eta\wedge\beta)|_S$. No factor has been moved across $\theta$. This explicit calculation fixes the convention later needed by \planned{residue_mul_map}.
\end{proof}

\section{The global pole-order comparison}

This section proves the quasi-isomorphism without local coordinates and without local-to-global cohomology arguments. The smooth-pair assumptions continue to hold. From the contraction step onwards assume that every positive integer is invertible in \(k\).

\begin{lemma}[Cartan calculus on the de Rham algebra]\label{lem:dr-cartan}
\lean{OrlikSolomon.DeRham.cartanFormula}
For a $k$-derivation $v:A\to A$, contraction $\iota_v$ is defined by the functional $\Om_A^1\to A$ taking $da$ to $v(a)$. It is an $A$-linear signed derivation of degree $-1$. The operator
\[
 \mathcal L_v=d\iota_v+\iota_vd
\]
is a $k$-linear derivation of degree zero, with $\mathcal L_v(a)=v(a)$ and $\mathcal L_v(da)=d(v(a))$. These constructions extend to localizations, and $\iota_{av}=a\iota_v$.
\uses{lem:dr-contraction,thm:dr-dga,lem:dr-localization}
\end{lemma}
\begin{proof}
Contraction and its signed product formula are Lemma~\ref{lem:dr-contraction}. Expand $d\iota_v(\alpha\beta)+\iota_vd(\alpha\beta)$ using the two signed Leibniz formulas. For $\deg\alpha=p$, the terms $\iota_v\alpha\,d\beta$ have coefficients $(-1)^{p-1}$ and $(-1)^p$, and the terms $d\alpha\,\iota_v\beta$ have coefficients $(-1)^p$ and $(-1)^{p+1}$. They cancel, leaving
$\mathcal L_v(\alpha)\beta+\alpha\mathcal L_v(\beta)$.
The values on $a$ and $da$ follow from the definitions and $d^2=0$. A derivation extends to localization by Lemma~\ref{lem:dr-localization}; use its functional on localized differentials to define contraction there. The scalar identity follows by comparing the functionals on one-forms and the word formula for contraction.
\end{proof}

\begin{lemma}[A logarithmic Euler contraction]\label{lem:dr-euler}
\lean{OrlikSolomon.DeRham.logEulerContraction}
Choose $v,h$ with $v(f)=1+fh$ as in Lemma~\ref{lem:dr-normal-derivation}, and set $E=fv$. On $\LL^\bullet$ the operator $\iota_E$ takes values in regular forms, preserves $f\LL^\bullet$, and descends to a degree $-1$ map $J$ on
\[
 V^\bullet=\LL^\bullet/f\LL^\bullet.
\]
Let $\bar d$ be its differential and $W$ left wedge multiplication by the class of $\theta$. Then
\[
 \bar dJ+J\bar d=0,\qquad WJ+JW=1.
\]
\uses{lem:dr-normal-derivation,lem:dr-cartan,lem:dr-log-dga}
\end{lemma}
\begin{proof}
For a regular form $\omega$, $\iota_E\omega=f\iota_v\omega$ is regular. For a log term,
\[
 \iota_E(\theta\wedge\eta)
 =\frac{E(f)}f\,\eta-\theta\wedge f\iota_v\eta
 =v(f)\eta-df\wedge\iota_v\eta
\]
is also regular. $A$-linearity implies $\iota_E(f\LL)\subseteq f\LL$, proving descent.

We claim $\mathcal L_E(\LL)\subseteq f\LL$. The log algebra is generated by the elements $a$, $da$ ($a\in A$), and $\theta$. On these generators,
\[
 \mathcal L_E(a)=f v(a),\qquad
 \mathcal L_E(da)=df\,v(a)+f\,d(v(a)),\qquad
 \mathcal L_E(\theta)=d(v(f))=h\,df+f\,dh.
\]
All are in $f\LL$ since $df=f\theta$. The degree-zero derivation rule and the fact that $f\LL$ is an ideal prove the claim on every log form. Cartan's formula now gives the first identity on $V$.

For the second identity, the contraction product rule gives
$\iota_E(\theta\wedge\xi)+\theta\wedge\iota_E\xi=v(f)\xi$.
Modulo $f\LL$, the scalar $v(f)=1+fh$ is $1$. Hence $JW+WJ=1$. This proof only needs the congruence modulo $f$; it does not assert that $v(f)=1$ in $A$.
\end{proof}

\begin{definition}[Pole filtration by subcomplexes]\label{def:dr-pole-filtration}
\lean{OrlikSolomon.DeRham.poleFiltration}
For $m\in\N$ put
\[
 F_m^n=f^{-m}\LL^n\subseteq\Om_{A_f}^n.
\]
Then $F_0=\LL$, the filtration is increasing and exhaustive, and every $F_m$ is a subcomplex.
\uses{lem:dr-log-dga,lem:dr-localization}
\end{definition}
\begin{proof}
The containment $f\LL\subseteq\LL$ proves $F_{m-1}\subseteq F_m$. Every localized form is $f^{-m}\omega$ for some regular $\omega$, hence lies in some $F_m$. Finally
\[
 d(f^{-m}\xi)=f^{-m}(d\xi-m\theta\wedge\xi)\in F_m.
\]
The use of log forms is essential here: the unadjusted modules $f^{-m}\Om_A^\bullet$ need not be subcomplexes.
\end{proof}

\begin{lemma}[The differential on a pole-order quotient]\label{lem:dr-pole-quotient}
\lean{OrlikSolomon.DeRham.poleQuotientDifferential}
For $m\geq1$ there is a graded $k$-linear isomorphism
\[
 V^\bullet=\LL^\bullet/f\LL^\bullet
 \xrightarrow{\sim} F_m^\bullet/F_{m-1}^\bullet,\qquad
 \bar\xi\longmapsto\overline{f^{-m}\xi}.
\]
Under it the differential on the right becomes
\[
 D_m=\bar d-mW
\]
on $V$.
\uses{def:dr-pole-filtration,lem:dr-log-dga}
\end{lemma}
\begin{proof}
Surjectivity follows from the definition of $F_m$. If $f^{-m}\xi=f^{-(m-1)}\eta$ for some $\eta\in\LL$, equality in the ambient localized module gives $\xi=f\eta$. This proves the kernel assertion and does not require choosing representatives in a localization quotient. The differential formula is the preceding displayed formula reduced modulo $f\LL$. Directly,
\[
 (d-m\theta\wedge)(f\eta)
 =f\bigl(d\eta+(1-m)\theta\wedge\eta\bigr),
\]
so it indeed descends. This also makes its well-definedness independent of the isomorphism argument.
\end{proof}

\begin{theorem}[Every positive pole-order quotient is contractible]\label{thm:dr-pole-contraction}
\lean{OrlikSolomon.DeRham.positivePoleQuotientContractible}
Assume $m\geq1$ is invertible in $k$. A contracting homotopy for $(V^\bullet,D_m)$, and hence for $F_m/F_{m-1}$, is
\[
 h_m=-\frac1m J.
\]
In particular $D_mh_m+h_mD_m=1$.
\uses{lem:dr-euler,lem:dr-pole-quotient}
\end{theorem}
\begin{proof}
Compute with the two identities of Lemma~\ref{lem:dr-euler}:
\[
 \begin{aligned}
 D_mh_m+h_mD_m
 &=-\frac1m\bigl(\bar dJ+J\bar d-m(WJ+JW)\bigr)\\
 &=-\frac1m(0-m\cdot1)=1.
 \end{aligned}
\]
This is the second use of rational scalars. The displayed equality is the desired Lean homotopy equation in each degree, not merely an assertion of vanishing cohomology. Contraction acts on representatives by the explicit word formula already proved.
\end{proof}

\begin{lemma}[Finite reduction through an acyclic quotient]\label{lem:dr-filtration-reduction}
\lean{OrlikSolomon.DeRham.reduceCocycleAndPrimitive}
Suppose $F_{m-1}\subseteq F_m$ is an inclusion of complexes with contractible quotient. Every cocycle in $F_m^n$ is cohomologous in $F_m$ to one in $F_{m-1}^n$. If a cocycle in $F_{m-1}^n$ is a boundary in $F_m$, it is already a boundary in $F_{m-1}$.
\end{lemma}
\begin{proof}
Let $z\in F_m^n$ be closed. Its quotient class $\bar z$ is closed, so $\bar z=d h\bar z$. Lift $h\bar z$ to $b\in F_m^{n-1}$. Then $z-db$ is a closed element of $F_{m-1}$. For the second assertion, write $z=dy$ in $F_m$ with $z\in F_{m-1}$. Then $\bar y$ is closed in the quotient and equals $d h\bar y$. Lift $h\bar y$ to $a\in F_m^{n-2}$. The element $y-da$ lies in $F_{m-1}$ and has differential $z$. Negative degrees cause no exception: they are zero, and all the displayed identities still make sense.
\end{proof}

\begin{theorem}[Global logarithmic comparison]\label{thm:dr-comparison}
\lean{OrlikSolomon.DeRham.logToLocalizationQuasiIso}
If $k$ contains $\mathbb Q$ and $(A,f)$ is a smooth principal pair, the inclusion
\[
 \LL^\bullet\hookrightarrow\Om_{A_f}^\bullet
\]
induces an isomorphism in every cohomological degree. It induces an algebra equivalence on total cohomology.
\uses{thm:dr-pole-contraction,lem:dr-filtration-reduction,def:dr-pole-filtration,def:dr-cohomology}
\end{theorem}
\begin{proof}
A localized cocycle belongs to some $F_m$ by exhaustivity. Apply Lemma~\ref{lem:dr-filtration-reduction} successively for $m,m-1,\ldots,1$ to obtain a cohomologous cocycle in $F_0=\LL$. This proves surjectivity. If a closed log form is the differential of a localized form $y$, put $y$ in some $F_m$. The second part of the same lemma successively replaces its primitive by one in $F_{m-1}$, until it is in $F_0$. This proves injectivity. Both arguments are finite inductions on the natural-number pole bound, so no spectral-sequence convergence or passage of cohomology through an infinite limit is required.

Inclusion preserves multiplication and units, so its bijective cohomology map is an algebra equivalence. The auxiliary normal derivation was used only to prove bijectivity; the cohomology map and its inverse are canonical. No localization-and-gluing theorem has been invoked. In particular we have not treated de Rham cohomology as an $A$-module, which it generally is not.
\end{proof}

\section{Cohomological residue, Gysin, and signs}

\begin{definition}[Untwisted cohomological maps]\label{def:dr-gysin}
\lean{OrlikSolomon.DeRham.cohomologicalResidueAndGysin}
Use the exact sequence in Theorem~\ref{thm:dr-residue-sequence} and the comparison isomorphism to define
\[
 r^n:H^n_{\dR}(A_f)\longrightarrow H^{n-1}_{\dR}(S),
 \qquad
 G^q:H^q_{\dR}(S)\longrightarrow H^{q+2}_{\dR}(A).
\]
The first is $H(\rho)$ after the inverse of logarithmic comparison. The second is the connecting map, with convention
\[
 G^q([\gamma])=[d\lambda],
 \quad
 \lambda\in\LL^{q+1},\quad \rho\lambda=\gamma,\quad d\gamma=0.
\]
Here $d\lambda$ is viewed as a regular form via exactness.
\uses{thm:dr-residue-sequence,thm:dr-comparison}
\end{definition}
\begin{proof}
The prompt permits the standard long-exact-sequence construction to be reused. We spell out its representative convention to fix its sign. Since $\rho d\lambda=-d\gamma=0$, exactness makes $d\lambda$ regular, and it is closed by $d^2=0$. Replacing $\lambda$ by another lift changes it by a regular form, so $d\lambda$ changes by a regular boundary. If $\gamma$ changes to $\gamma+d\eta$, choose a log lift $\mu$ of $\eta$ and replace $\lambda$ by $\lambda-d\mu$; the anticommutation $\rho d=-d\rho$ gives the correct new residue, and $d\lambda$ is unchanged. Combining these two observations proves that the displayed convention gives a well-defined map on ordinary cohomology of $S$, including its identification with cohomology of the shifted target. The residue is well defined because comparison is an isomorphism.
\end{proof}

\begin{lemma}[Representative formulas and exactness]\label{lem:dr-exactness}
\lean{OrlikSolomon.DeRham.gysinResidueExactness}
The cohomological sequence is
\[
 \cdots\to H^{n-2}_{\dR}(S)\xrightarrow{G}H^n_{\dR}(A)
 \xrightarrow{j_*}H^n_{\dR}(A_f)
 \xrightarrow{r}H^{n-1}_{\dR}(S)
 \xrightarrow{G}H^{n+1}_{\dR}(A)\to\cdots.
\]
For a closed log representative $\omega+\theta\wedge\eta$,
\[
 r([\omega+\theta\wedge\eta])=[\eta|_S].
\]
For a closed $\gamma\in\Om_S^q$ and any regular lift $\eta\in\Om_A^q$,
\[
 G([\gamma])=[-\theta\wedge d\eta]\in H^{q+2}_{\dR}(A).
\]
The expression inside the last brackets is regular.
\uses{def:dr-gysin,lem:dr-regularity,thm:dr-residue-sequence,thm:dr-comparison}
\end{lemma}
\begin{proof}
The long exact sequence of the short exact sequence has terms $H^n(\Om_S[-1])=H^{n-1}(\Om_S)$ and connecting degree $+1$, giving precisely the displayed degrees. Replace $H^n(\LL)$ by $H^n(A_f)$ using comparison. The residue formula is its definition. For Gysin use the lift $\lambda=\theta\wedge\eta$; its differential is $-\theta\wedge d\eta$. Since $(d\eta)|_S=d\gamma=0$, Lemma~\ref{lem:dr-regularity} proves regularity directly. These formulas introduce no additional $(-1)^q$.

The minus differential on the shifted complex makes $\rho$ a chain map. It does not negate the form normalization $\rho(\theta\wedge\eta)=\eta|_S$, nor does it insert a sign in the right product formula below. A library connecting-map convention must be reconciled with the explicit formula $[d\lambda]$ when doing the Lean transport.
\end{proof}

\begin{theorem}[Cohomological right projection formula]\label{thm:dr-projection}
\lean{OrlikSolomon.DeRham.residueRightProduct}
For $a\in H^n_{\dR}(A_f)$ and $b\in H^q_{\dR}(A)$,
\[
 r(a\,j_*b)=r(a)\,i_*b.
\]
In particular this has the exact sign required when $n=1$.
\uses{lem:dr-form-projection,thm:dr-comparison,def:dr-gysin}
\end{theorem}
\begin{proof}
Choose a closed log form $\alpha$ representing $a$, which exists by logarithmic comparison, and a closed regular form $\beta$ representing $b$. Their product is a closed log form representing $a\,j_*b$. On representatives,
\[
 \rho\bigl((\omega+\theta\wedge\eta)\wedge\beta\bigr)
 =(\eta\wedge\beta)|_S=(\eta|_S)\wedge(\beta|_S).
\]
Passing to cohomology proves the assertion. The inverse comparison respects products because comparison itself is a bijective algebra map. No appeal to a signed abstract module convention is needed: the right-factor computation above supplies the sign explicitly.
\end{proof}

\section{Naturality for associated equations}

\begin{theorem}[Naturality of the log residue sequence]\label{thm:dr-log-naturality}
\lean{OrlikSolomon.DeRham.logSequenceNaturality}
Let $(A,f)$ and $(A',f')$ be smooth principal pairs and let $\phi:A\to A'$ satisfy $\phi(f)=u f'$ for a unit $u\in(A')^\times$. Let $\psi:S\to S'$ and $\chi:A_f\to A'_{f'}$ be the induced quotient and localization maps. Then $\chi_*$ takes log forms to log forms, commutes with their differentials, and satisfies
\[
 \rho'\chi_*=\psi_*\rho.
\]
Together with $\phi_*$ this is a morphism of the short exact sequences of complexes. No flatness of $\phi$ is required.
\uses{thm:dr-dga,lem:dr-localization,def:dr-residue,thm:dr-residue-sequence}
\end{theorem}
\begin{proof}
The quotient map exists because $\phi(f)$ is zero modulo $f'$, and the localization map exists because it is invertible after inverting $f'$. Their compatibility with the original algebra maps follows on generators from the universal properties. Differentiate the equation in the target localization:
\[
 \chi_*\theta=\frac{d(uf')}{uf'}=\frac{df'}{f'}+\frac{du}{u}
 =\theta'+\frac{du}{u}.
\]
The second term is regular on $A'$ because $u$ is an ambient unit. Thus
\[
 \chi_*(\omega+\theta\wedge\eta)
 =\phi_*\omega+\frac{du}{u}\wedge\phi_*\eta
   +\theta'\wedge\phi_*\eta
\]
is logarithmic, and its residue is $(\phi_*\eta)|_{S'}=\psi_*(\eta|_S)$. Commutation with differentials holds by functoriality of forms; the other square is simply restriction of the localization square. The proof neither tensors an exact sequence along $\phi$ nor assumes its flatness. An arbitrary fixed morphism of divisor pairs in the prompt has exactly these maps by their stated commuting squares and the universal properties.
\end{proof}

\begin{theorem}[Naturality of residue and Gysin]\label{thm:dr-naturality}
\lean{OrlikSolomon.DeRham.residueAndGysinNaturality}
Under the hypotheses above,
\[
 r'\,\chi_*=\psi_*\,r,\qquad G'\,\psi_*=\phi_*\,G.
\]
The constructions are unchanged by replacing $f$ by an associated equation in the same ambient algebra.
\uses{thm:dr-log-naturality,thm:dr-comparison,def:dr-gysin}
\end{theorem}
\begin{proof}
For residue, the square of log inclusions into localized forms commutes. Its two cohomology maps are isomorphisms, so their inverses also commute with the vertical maps: compose the desired equality with one of the isomorphisms and use its injectivity. Combine this with the form-residue square.

For Gysin, let $\lambda$ lift a closed representative $\gamma$ of a source divisor class. The target log form $\chi_*\lambda$ lifts $\psi_*\gamma$ by Theorem~\ref{thm:dr-log-naturality}. Its differential is $\chi_*d\lambda$, and since $d\lambda$ is regular this is the image under $\phi_*$ of that regular form. Thus the representative convention gives $G'\psi_*[\gamma]=\phi_*G[\gamma]$, with no sign. For an associated equation apply the same argument with $\phi$ the identity. The localized and log algebras agree as subalgebras, since the extra $du/u$ is regular. The choices of normal derivation used to prove comparison play no role in these natural maps.
\end{proof}

\section{Construction of the fixed ordinary-theory instance}

\begin{definition}[The Laurent-twisted functor]\label{def:dr-twisted-functor}
\lean{OrlikSolomon.DeRham.twistedCohomologyFunctor}
Let $k$ be a field of characteristic zero and $B=k[t,t^{-1}]$. For each smooth $k$-algebra put
\[
 C(A)=H^\bullet_{\dR}(A/k)\otimes_k B,\qquad
 C(A)_{p,w}=H^p_{\dR}(A/k)t^w,
 \qquad \mathcal B_w=kt^w.
\]
Use zero for $p<0$. Define $C(\phi)=H^\bullet(\phi_*)\otimes1$. The $B$-algebra structure uses the coefficient factor, and the $k$-algebra structure is its restriction.
\uses{def:dr-cohomology}
\end{definition}
\begin{proof}
The Laurent monomials form a $k$-basis indexed by $\Z$. Tensoring a vector space with $B$ is therefore a direct sum of copies indexed by the weight, not a product. Decomposing total cohomology by degree gives the internal $(p,w)$ grading. The product adds both degrees and the unit has degree $(0,0)$. Coefficients are central because their cohomological degree is zero. Functoriality, the scalar tower, and the stated coefficient embedding follow directly from tensor products and the functor on cohomology. As allowed by the prompt, implementers may package these standard direct-sum and \planned{GradedAlgebra} constructions without reproving their general infrastructure here.
\end{proof}

\begin{lemma}[Point, grading, and affine-line fields]\label{lem:dr-ordinary-basic}
\lean{OrlikSolomon.DeRham.ordinaryTheoryBasicFields}
The functor and grading above satisfy all fields of the fixed interface preceding the Gysin and residue fields. In particular the point map $B\to C(k)$ is bijective, nonzero cohomological point degrees vanish, and $C(A)\to C(A[z])$ is bijective for every smooth $A$.
\uses{def:dr-twisted-functor,thm:dr-homotopy,thm:dr-dga}
\end{lemma}
\begin{proof}
The coefficient grading has $\mathcal B_0=k\cdot1$, so its structural map from $k$ is bijective. The dg algebra functor preserves form degree, hence cohomology degree, and tensoring by $B$ preserves weight. Graded commutativity of forms descends to the cohomology algebra with sign $(-1)^{pq}$; this is the integer \planned{negOnePow} sign in the interface, after extending by zero in negative degrees.

Every $k$-derivation $k\to M$ is zero, so $\Om^1_{k/k}=0$. Thus the point complex is $k$ in degree zero and zero elsewhere, which gives $C(k)=B$ and the required point fields. Since $k$ is a characteristic-zero field, the scalar image of each positive integer is nonzero and therefore invertible. Theorem~\ref{thm:dr-homotopy} applies to every $A$, and tensoring its actual inverse with $B$ proves bijectivity for \planned{polynomialC}. This is where the field and characteristic hypotheses supply the rational scalars used in the homotopy.
\end{proof}

\begin{definition}[Twisted Gysin and residue maps]\label{def:dr-twisted-maps}
\lean{OrlikSolomon.DeRham.twistedGysinAndResidue}
For a fixed \planned{DivisorPair}, transport the canonical quotient/localization construction along its specified maps. Extend $G$ and $r$ $B$-linearly and set
\[
 g(y\,t^w)=G(y)t^{w+1},\qquad
 R(x\,t^w)=r(x)t^{w-1}.
\]
These are the proposed fields \planned{gysin} and \planned{residue}.
\uses{def:dr-gysin,def:dr-twisted-functor}
\end{definition}
\begin{proof}
Multiplication by $t$ and by $t^{-1}$ is $B$-linear, so the displayed formulas extend to well-defined $B$-linear maps. The quotient and localization isomorphisms are characterized by commutation with $i,j$; the given kernel, surjectivity, and \planned{IsLocalization.Away} hypotheses provide them. Transporting the construction across those canonical isomorphisms is the standard task explicitly assigned to the implementers in the prompt.
\end{proof}

\begin{theorem}[All localization and compatibility fields]\label{thm:dr-ordinary-localization}
\lean{OrlikSolomon.DeRham.ordinaryTheoryLocalizationFields}
The maps $g,R$ have shifts $(2,1)$ and $(-1,-1)$, respectively. They satisfy
\[
 \im g=\ker C(j),\qquad
 \im C(j)=\ker R,\qquad
 \im R=\ker g,
\]
commute with every fixed morphism of divisor pairs, and satisfy
\[
 R\bigl(\alpha\,C(j)\beta\bigr)=R(\alpha)\,C(i)\beta
 \quad\text{for }\alpha\in C(A_f)_{1,1},\ \beta\in C(A).
\]
\uses{def:dr-twisted-maps,lem:dr-exactness,thm:dr-projection,thm:dr-naturality}
\end{theorem}
\begin{proof}
The cohomological degrees of $G$ and $r$, together with the explicit powers of $t$, give the two shifts. In each bidegree, the three kernel/image identities are precisely the corresponding portions of Lemma~\ref{lem:dr-exactness}, with indices shifted. Exactness of the total direct sums follows by lifting the finitely many homogeneous components of an element. Equivalently, expand in Laurent monomials and use degreewise exactness. This does not require a field-dimension argument.

A morphism of pairs preserves the Laurent factor, so both naturality identities follow from Theorem~\ref{thm:dr-naturality}. For the last field first take $\alpha=a t$ and $\beta=b t^v$ homogeneous. Theorem~\ref{thm:dr-projection} gives the entire sign calculation:
\[
 \begin{aligned}
 R\bigl((a t)\,C(j)(b t^v)\bigr)
 &=r(a\,j_*b)t^{v}\\
 &=r(a)i_*b\,t^v
 =R(a t)\,C(i)(b t^v).
 \end{aligned}
\]
For arbitrary $\beta$, decompose it into finitely many homogeneous terms and add the identities. The condition on $\alpha$ identifies it with $a t$; no extra degree sign occurs. This proves exactly \planned{residue_mul_map}, with $\alpha$ first.
\end{proof}

\begin{theorem}[The nonzero de Rham instance]\label{thm:dr-instance}
\lean{OrlikSolomon.DeRham.ordinaryTheory}
For every characteristic-zero field $k$, the preceding data define a value of the fixed
\[
 \planned{OrdinaryTheory}\ k\ \mathcal B
 \quad\text{with }R=k,\quad B=k[t,t^{-1}].
\]
The coefficient ring and the cohomology algebra of the point are nonzero. All fields hold for every smooth algebra and every divisor pair of the prompt, without arrangements, connectedness, nonemptiness, or flatness assumptions on morphisms.
\uses{lem:dr-ordinary-basic,thm:dr-ordinary-localization,def:cohomology}
\end{theorem}
\begin{proof}
Use Definition~\ref{def:dr-twisted-functor} for the functor and grading, Lemma~\ref{lem:dr-ordinary-basic} for the point, grading, and affine-line fields, Definition~\ref{def:dr-twisted-maps} for the two maps, and Theorem~\ref{thm:dr-ordinary-localization} for all remaining fields. A divisor pair has exactly the hypotheses used above: ambient formal smoothness, quotient formal smoothness, and a non-zero-divisor equation. Its supplied isomorphisms place the construction in the fixed types.

Evaluation $t\mapsto1$ is a $k$-algebra map $B\to k$ splitting the coefficient inclusion. Since a field is nonzero, $B$ is nonzero. The point identification then makes $C(k)$ nonzero. No assertion that $C(A)$ is nonzero for every object is required or intended: the zero algebra is a valid object and has zero cohomology.
\end{proof}

\paragraph{Field-by-field checklist.}
Every entry below is a field or an instance field of the fixed structure. Definitions are data constructors; their cited proofs establish the required laws.
{\small
\begin{longtable}{@{}p{0.53\textwidth}p{0.42\textwidth}@{}}
\toprule
Fixed field & Proven by\\\midrule\endhead
\nolinkurl{bijective_algebraMap_gradeZero} & Lemma~\ref{lem:dr-ordinary-basic}.\\
\nolinkurl{C} & Definition~\ref{def:dr-twisted-functor}; Theorem~\ref{thm:dr-dga}.\\
\nolinkurl{algebra} & Definition~\ref{def:dr-twisted-functor}.\\
\nolinkurl{isScalarTower} & Definition~\ref{def:dr-twisted-functor}.\\
\nolinkurl{grading} & Definition~\ref{def:dr-twisted-functor}.\\
\nolinkurl{gradedAlgebra} & Definition~\ref{def:dr-twisted-functor}.\\
\nolinkurl{algebraMap_mem_grading} & Definition~\ref{def:dr-twisted-functor}; Lemma~\ref{lem:dr-ordinary-basic}.\\
\nolinkurl{isHomogeneousOfShift_map} & Lemma~\ref{lem:dr-ordinary-basic}.\\
\nolinkurl{isGradedCommutative} & Lemma~\ref{lem:dr-ordinary-basic}.\\
\nolinkurl{bijective_algebraMap_point} & Lemma~\ref{lem:dr-ordinary-basic}.\\
\nolinkurl{eq_zero_of_mem_grading_point} & Lemma~\ref{lem:dr-ordinary-basic}.\\
\nolinkurl{bijective_map_polynomialC} & Theorem~\ref{thm:dr-homotopy}; Lemma~\ref{lem:dr-ordinary-basic}.\\
\nolinkurl{gysin} & Definition~\ref{def:dr-twisted-maps}.\\
\nolinkurl{residue} & Definition~\ref{def:dr-twisted-maps}.\\
\nolinkurl{isHomogeneousOfShift_gysin} & Theorem~\ref{thm:dr-ordinary-localization}.\\
\nolinkurl{isHomogeneousOfShift_residue} & Theorem~\ref{thm:dr-ordinary-localization}.\\
\nolinkurl{exact_gysin_map} & Theorem~\ref{thm:dr-ordinary-localization}.\\
\nolinkurl{exact_map_residue} & Theorem~\ref{thm:dr-ordinary-localization}.\\
\nolinkurl{exact_residue_gysin} & Theorem~\ref{thm:dr-ordinary-localization}.\\
\nolinkurl{gysin_map} & Theorem~\ref{thm:dr-naturality}; Theorem~\ref{thm:dr-ordinary-localization}.\\
\nolinkurl{residue_map} & Theorem~\ref{thm:dr-naturality}; Theorem~\ref{thm:dr-ordinary-localization}.\\
\nolinkurl{residue_mul_map} & Theorem~\ref{thm:dr-projection}; its twisted computation in Theorem~\ref{thm:dr-ordinary-localization}.\\
\bottomrule
\end{longtable}
}

\begin{lemma}[The normalized logarithmic class]\label{lem:dr-tau}
\lean{OrlikSolomon.DeRham.logarithmicClass_eq_dlog}
For the standard pair $(k[x],x)$, the normalized class $\tau$ of the first blueprint is $[dx/x]\,t\in C(k[x,x^{-1}])_{1,1}$. For an invertible regular function $a$ on any smooth algebra, its abstract logarithmic class is $[da/a]\,t$.
\uses{thm:dr-instance,def:dr-residue,lem:gm,def:log}
\end{lemma}
\begin{proof}
The form $dx/x$ is closed and has form residue $1$, so the twisted residue of $[dx/x]t$ is $1$ in degree $(0,0)$. The uniqueness assertion defining the normalized class in the first blueprint identifies it with $\tau$. For any algebra map $k[x,x^{-1}]\to A$ taking $x$ to a unit $a$, functoriality takes $dx/x$ to $da/a$. This gives the second assertion directly, without reproving the abstract logarithm identities.
\end{proof}

\begin{corollary}[Affine arrangement corollary]\label{cor:dr-affine}
\lean{OrlikSolomon.DeRham.affineOrlikSolomonEquiv}
For every split affine hyperplane arrangement $\cA$ over a characteristic-zero field $k$,
\[
 \OS_k(\cA)\simeq_k\bigoplus_{q\geq0}H^q_{\dR}(U_{\cA}/k),
 \qquad e_H\longmapsto[df_H/f_H].
\]
\uses{thm:dr-instance,lem:dr-tau,thm:affine}
\end{corollary}
\begin{proof}
Instantiate the existing affine structure theorem, \extlabel{thm:affine}, with Theorem~\ref{thm:dr-instance}. Its diagonal cohomology algebra is $\bigoplus_q H^q_{\dR}(U)t^q$. Sending $x t^q$ to $x$ on this diagonal is an algebra equivalence: its inverse sends a homogeneous degree-$q$ class to $x t^q$, and multiplication adds degrees. Lemma~\ref{lem:dr-tau} gives the formula on generators. Thus the conclusion no longer has a hypothetical cohomology-theory parameter. No part of the arrangement proof is repeated.
\end{proof}

\begin{corollary}[Projective arrangement corollary]\label{cor:dr-projective}
\lean{OrlikSolomon.DeRham.projectiveOrlikSolomonEquiv}
For a nonempty split projective arrangement $\cP$ in $\PP^n_k$,
\[
 \pOS_k(\cP)\simeq_k\bigoplus_{q\geq0}H^q_{\dR}(U_{\cP}/k),
 \qquad e_i-e_j\longmapsto
 \left[\frac{d(\ell_i/\ell_j)}{\ell_i/\ell_j}\right].
\]
The map is independent of all equations and of the chosen hyperplane at infinity.
\uses{thm:dr-instance,lem:dr-tau,thm:projective}
\end{corollary}
\begin{proof}
Instantiate the existing projective theorem, \extlabel{thm:projective}, and remove the diagonal powers of $t$ as in the preceding corollary. The nonempty projective complement is affine by deconing, so its cohomology here is the de Rham complex of its smooth coordinate algebra. The abstract ratio generator pulls back to the displayed form by Lemma~\ref{lem:dr-tau}. Independence is part of the already-proved projective theorem. The empty projective arrangement requires its separate projective calculation and is not included in this corollary.
\end{proof}

\section{Worked regressions and excluded cases}

Unless stated otherwise $k$ is a characteristic-zero field. These calculations test the maps and signs, not just the dimensions of the answer.

\begin{proposition}[The punctured affine line]\label{prop:dr-gm-example}
\lean{OrlikSolomon.DeRham.puncturedLineRegression}
For $A=k[x]$ and $f=x$,
\[
 H^0_{\dR}(A_f)=k,\qquad H^1_{\dR}(A_f)=k[dx/x],
 \qquad H^n_{\dR}(A_f)=0\ (n>1),
\]
and $r([dx/x])=1$. The twisted normalized class is $[dx/x]t$.
\uses{lem:dr-localization,def:dr-residue,def:dr-twisted-maps}
\end{proposition}
\begin{proof}
The complex is $k[x,x^{-1}]\xrightarrow{d}k[x,x^{-1}]dx$ with $d(x^m)=m x^{m-1}dx$ for all $m\in\Z$. Only the constant monomial has zero derivative. Every one-form monomial $x^r dx$ except $r=-1$ has primitive $x^{r+1}/(r+1)$; the coefficient of $x^{-1}dx$ in any derivative is zero. Thus it spans the cokernel and is nonzero. Higher forms vanish because the module of one-forms has rank one. Form residue sends $dx/x$ to $1$, and the twist formula gives the final assertion.
\end{proof}

\begin{proposition}[A unit equation]\label{prop:dr-unit-example}
\lean{OrlikSolomon.DeRham.unitEquationRegression}
If $f$ is a unit in a smooth $A$, then $S=0$, $A_f\simeq A$, and $\LL^\bullet=\Om_A^\bullet$. Both residue and Gysin are zero, and all localization exactness fields reduce to the identity isomorphism on $C(A)$.
\uses{def:dr-log,def:dr-residue,def:dr-gysin}
\end{proposition}
\begin{proof}
The quotient by a unit is zero and localization at a unit changes no algebra. Since $df/f$ is already regular, adjoining its wedges changes no forms. The zero algebra has zero forms in every degree, so the residue sequence is $0\to\Om_A\to\Om_A\to0\to0$. Its cohomological maps have the asserted values. In the global contraction proof one may take $v=0$ and $h=-f^{-1}$; every pole-order quotient is zero. This confirms that no step assumed a nonempty divisor.
\end{proof}

\begin{proposition}[A divisor on one connected component]\label{prop:dr-product-example}
\lean{OrlikSolomon.DeRham.disconnectedPairRegression}
For $A=k\times k[z]$ and $f=(1,z)$,
\[
 S\simeq k,\qquad A_f\simeq k\times k[z,z^{-1}],\qquad
 H^0_{\dR}(A_f)=k\times k,\qquad
 H^1_{\dR}(A_f)=k[(0,dz/z)].
\]
All higher groups vanish, and residue takes $(0,dz/z)$ to $1$ on $S$.
\uses{prop:dr-gm-example,def:dr-residue,lem:dr-normal-derivation}
\end{proposition}
\begin{proof}
Differentials split over a finite product. To see this directly, for an idempotent $e$ the Leibniz rule gives $(1-2e)de=0$, and $(1-2e)^2=1$, so $de=0$. A derivation into a product-ring module therefore preserves its two idempotent summands; its universal differential module and all exterior powers split accordingly. The de Rham differential respects the splitting, so the complex is a product of the two complexes, and the product is finite.

The quotient by $(1,z)$ kills the first component and evaluates the second at zero. Localization keeps the first component and inverts $z$ in the second. Apply the punctured-line calculation and the point calculation componentwise. The logarithmic form is $(0,dz/z)$, so its residue is $1$ in the second divisor component, which is all of $S$. Gysin is zero here because its target groups have degree at least two.

For the global normal derivation take $v=(0,\partial_z)$. Then $v(f)=(0,1)=1+f(-1,0)$. The form $df=(0,dz)$ is not unimodular on all of $A$: every $A$-linear functional on it has first component zero. It is unimodular modulo $f$, exactly as required by the proof. This is the example denoted $k\times k[t]$ in the prompt, with its geometric variable renamed to avoid confusion with the coefficient variable.
\end{proof}

\begin{proposition}[The zero algebra]\label{prop:dr-zero-example}
\lean{OrlikSolomon.DeRham.zeroAlgebraRegression}
For $A=0$, every form and cohomology module is zero, $C(A)$ is the zero $B$-algebra, and the divisor-pair fields hold for its unique equation.
\uses{def:dr-forms,def:dr-twisted-functor,prop:dr-unit-example}
\end{proposition}
\begin{proof}
A unital module over the zero ring is zero because $m=1m=0m=0$. Thus $\Om_A^n=0$ including degree zero. The unique element of the zero ring is a unit and is a non-zero-divisor in the injectivity sense: the condition $ax=0\Rightarrow x=0$ is vacuous except for $x=0$. Its quotient and localization are again zero, so the preceding unit case applies. The zero $k$-algebra is formally smooth: if a nilpotent quotient of a ring is zero, the whole ring is zero, so any required lift from it exists. It is finitely presented as $k/(1)$. Therefore it is genuinely among the objects quantified over by the fixed interface. No nontriviality instance may be imposed on general \planned{SmoothAlg} objects.
\end{proof}

\begin{proposition}[Why the crossing divisor is excluded]\label{prop:dr-crossing-example}
\lean{OrlikSolomon.DeRham.singularDivisorCounterexample}
For $A=k[x,y]$ and $f=xy$, $f$ is a non-zero-divisor but $S=A/(xy)$ is not smooth. The regularity criterion needed for the residue definition fails: with $\eta=y\,dx$,
\[
 \frac{d(xy)}{xy}\wedge\eta=dy\wedge dx
 \quad\text{is regular},\qquad
 \eta|_S\ne0\text{ in }\Om_S^1.
\]
\uses{lem:dr-quotient,lem:dr-conormal}
\end{proposition}
\begin{proof}
The polynomial algebra is a domain, so $xy$ is a non-zero-divisor. The conormal vector is $e=y\,dx+x\,dy$. Modulo the maximal ideal $(x,y)$ its image in $\Om_A^1\otimes_A k$ is zero. A functional taking $e$ to $1$ would remain such after this reduction, which is impossible. If $S$ were formally smooth, Lemma~\ref{lem:dr-conormal} would provide that functional. Thus it is not smooth.

In the localization, $d(xy)/(xy)=dx/x+dy/y$, which gives the displayed regular wedge. By Lemma~\ref{lem:dr-quotient},
\[
 \Om_S^1=(S\,dx\oplus S\,dy)/S(y\,dx+x\,dy).
\]
If $y\,dx$ were zero, there would be $a\in S$ with $ay=y$ and $ax=0$. The annihilator of $x$ in $S$ is $(y)$: lifting $ax=0$ to $k[x,y]$ and cancelling $x$ proves this. Hence $a=yb$ for some $b$, and $y^2b=y$. Reducing modulo $x$ gives this equality in $k[y]$, impossible by its coefficient of $y$. Thus $\eta|_S\ne0$. The same log form has both a regular presentation and a logarithmic presentation whose proposed residues disagree. This identifies the exact failure, not merely a missing smoothness hypothesis.
\end{proof}

\begin{proposition}[Unit changes, a non-flat map, and ramification]\label{prop:dr-morphism-examples}
\lean{OrlikSolomon.DeRham.morphismRegressions}
The following are tests of the precise morphism condition.
\begin{enumerate}
\item For $k[z]\to k[x,y,y^{-1}]$, $z\mapsto xy$, and target equation $f'=x$, the residue of $d(xy)/(xy)$ is $1$.
\item A genuinely non-flat morphism of admissible pairs is
\[
 (k[z,v],z)\longrightarrow(k[x,y,y^{-1}],x),\qquad z\mapsto xy,\quad v\mapsto0.
\]
Residue and Gysin naturality still hold.
\item The map $(k[z],z)\to(k[x],x)$ sending $z$ to $x^2$ is not a morphism of the fixed pairs. Its pulled-back logarithmic form has residue $2$, not $1$.
\end{enumerate}
\uses{thm:dr-naturality,def:dr-residue}
\end{proposition}
\begin{proof}
In the first case $u=y$ is an ambient unit and
$d(xy)/(xy)=dx/x+dy/y$. The second summand is regular in the ambient algebra, so residue is $1$. This map is actually flat: the target is isomorphic as a $k[z]$-algebra to $k[z,y,y^{-1}]$, with inverse coordinate change $x=z/y$. It is free as a $k[z]$-module on the Laurent monomials in $y$.

In the second case the ambient and divisor algebras on both sides are smooth, the equations are non-zero-divisors, and the image of $z$ is the unit $y$ times $x$. The quotient maps send $v$ to zero, and the localization maps exist by inverting $z$ and $x$. This gives all the fixed morphism fields. It is not flat: multiplication by $v$ is injective on $k[z,v]$, but after tensoring with the nonzero target algebra it is the zero endomorphism, which is not injective. Naturality is Theorem~\ref{thm:dr-naturality}, whose proof made no flatness assumption.

Finally $x^2$ is not associated to $x$ in $k[x]$, since their ratio $x$ is not a unit. The form pulls back to $2dx/x$, of residue $2$. The divisor map is the identity on $k$, so the proposed naturality equality without the associated-equation restriction would force $2=1$. This verifies the necessity of excluding the ramified map in the fixed interface.
\end{proof}

\begin{proposition}[The two-dimensional torus without K\"unneth]\label{prop:dr-torus-example}
\lean{OrlikSolomon.DeRham.twoTorusRegression}
For $T=k[x^{\pm1},y^{\pm1}]$, put $e_x=dx/x$ and $e_y=dy/y$. Then
\[
 H^0_{\dR}(T)=k,\qquad H^1_{\dR}(T)=k[e_x]\oplus k[e_y],\qquad
 H^2_{\dR}(T)=k[e_x\wedge e_y],
\]
with all higher groups zero and the exterior multiplication on these generators. The residue along $x=0$ in ambient $k[x,y^{\pm1}]$ sends $[e_x\wedge e_y]$ to $[e_y]$. Along $y=0$ it sends the same ordered class to $-[e_x]$.
\uses{thm:dr-dga,lem:dr-localization,lem:dr-contraction,def:dr-residue}
\end{proposition}
\begin{proof}
The forms have basis $x^a y^b e_I$ with $(a,b)\in\Z^2$ and $I$ an exterior word in $e_x,e_y$. The complex is the direct sum, over those exponent pairs, of the exterior algebra on a two-dimensional $k$-space with differential given by left wedge multiplication by $a e_x+b e_y$. For $(a,b)\ne(0,0)$ choose a functional taking this vector to $1$: use the value $1/a$ on $e_x$ and zero on $e_y$ when $a\ne0$, and otherwise use $1/b$ on $e_y$. Contraction by this functional is a contracting homotopy by Lemma~\ref{lem:dr-contraction}. At $(0,0)$ the differential is zero. Since every Laurent form has finitely many exponents, summing these component homotopies is well defined. Therefore exactly the zero-exponent exterior algebra remains in cohomology. This also proves the multiplicative assertion.

For $x=0$, $e_x$ is the left logarithmic factor and $e_y$ is regular, so residue of their product is $e_y$. For $y=0$, rewrite $e_x\wedge e_y=-e_y\wedge e_x$ and apply the same normalization. This is consistent with the right projection formula; no K\"unneth theorem has been used.
\end{proof}

\begin{proposition}[The obstruction in positive characteristic]\label{prop:dr-characteristic-p}
\lean{OrlikSolomon.DeRham.positiveCharacteristicFailure}
If $k$ is a field of characteristic $p>0$, the polynomial-invariance field fails for the proposed ordinary de Rham model. In $k[z]$, the form $z^{p-1}dz$ is a nonzero cohomology class, and $z^p$ is a nonconstant degree-zero cocycle.
\uses{thm:dr-dga}
\end{proposition}
\begin{proof}
Every one-form on $k[z]$ is closed. In the derivative of a polynomial, the coefficient of $z^{p-1}dz$ is $p$ times the coefficient of $z^p$, hence zero. Thus the displayed one-form is not exact. Also $d(z^p)=0$, so even degree-zero cohomology is larger than $k$. The map induced by $k\to k[z]$ cannot be bijective.

The failure also appears in the pole argument for $f=z$: the closed Laurent form $z^{-p-1}dz$ is not exact, since its only possible monomial primitive would involve division by $-p$. It cannot be represented by a log form, whose exponents in degree one are at least $-1$; the difference would still have a nonzero coefficient of $z^{-p-1}dz$ and could not be a derivative. Thus the positive pole quotient at $m=p$ cannot be contracted as above. Both uses of rational scalars are real requirements.
\end{proof}

\section{Hypothesis audit and formalization dependencies}

The proofs have been arranged so that each use of the hypotheses can be isolated in a Lean lemma. They establish the following precise division of work.
{\small
\begin{longtable}{@{}p{0.30\textwidth}p{0.65\textwidth}@{}}
\toprule
Hypothesis & Where it is used\\\midrule\endhead
Commutative base ring $k$ & Sufficient for construction and functoriality of the dg algebra, quotient/localization formulas, exterior calculus, and all form-level signs.\\
Formal smoothness of $A$ & Supplies projectivity of $\Om_A^1$, hence torsion-freeness of its exterior powers with respect to $f$, and the lift of a normal functional to an $A$-valued derivation.\\
Formal smoothness of $S$ & Supplies the section of $A/(f^2)\to S$, hence the split conormal line and the normal functional modulo $f$.\\
$f$ a non-zero-divisor & Identifies $(f)/(f^2)$ with $S$ and ensures regular forms inject into localized forms.\\
Positive integers invertible in $k$ & Used in the integration homotopy and in $h_m=-J/m$ for every $m\geq1$, and in the explicit characteristic-zero regressions.\\
$k$ a characteristic-zero field & Guarantees those invertibilities and nonzero coefficients. No separate geometric use of the field property occurs in the construction.\\
Finite presentation of smooth algebras & Required to belong to the fixed category. The algebraic arguments themselves only use the formal-smoothness properties just listed.\\
Associated equations for a pair morphism & Gives $\phi(f)=u f'$ with $du/u$ regular in the target ambient algebra, hence the exact residue normalization and naturality.\\
\bottomrule
\end{longtable}
}
In particular the proofs work over a nonzero commutative $\mathbb Q$-algebra base with the same formal-smoothness and regularity assumptions. The construction required by the prompt is its field case. One can weaken the smooth-pair input further to injectivity of regular forms into the localization and existence of a derivation $v(f)\equiv1\bmod f$: reducing that derivation modulo $f$ supplies the contraction needed in Lemma~\ref{lem:dr-modf-koszul}. This observation is a reading of the proofs, not a proposed weakening of the fixed interface.

\paragraph{Dependency graph.}
All mathematical nodes have a unique kind-prefixed \planned{dr-} label and a proposed declaration in \planned{OrlikSolomon.DeRham}. The \planned{uses} annotations give the direct dependencies. References to \planned{def:cohomology}, \planned{lem:gm}, \planned{def:log}, \planned{thm:affine}, and \planned{thm:projective} refer to the first blueprint and are external nodes. They are not restated or changed here. The following groups summarize the development order.

\begin{enumerate}
\item \textbf{Foundations.} Exterior base change and quotients; contraction; the polynomial de Rham dg algebra; descent through the canonical polynomial presentation; localization and quotient formulas; cohomology products. These use no smoothness or characteristic-zero assumptions.
\item \textbf{Two independent branches.} After the dg algebra is constructed, polynomial decomposition and integration establish affine-line invariance. Independently, smooth-pair algebra establishes the split conormal sequence, regular-form injectivity, and the global normal derivation. The contraction API can also be implemented independently of the de Rham differential.
\item \textbf{Log forms.} Define the actual submodules in localized forms, prove the regularity criterion, and construct the residue sequence with target differential $-d$. Do not add a pre-existing residue or comparison isomorphism as an assumption.
\item \textbf{Pole comparison.} Implement Cartan calculus, then the Euler identities on $\LL/f\LL$. Build $F_m$, identify its quotients, and formalize the single calculation $D_mh_m+h_mD_m=1$. Use the two elementary finite-reduction arguments to prove comparison.
\item \textbf{Maps and compatibilities.} Construct the cohomological maps using the existing long exact sequence. Naturality of log complexes under unit changes can be proved independently of pole comparison; naturality of localized cohomological residue uses comparison afterwards. Fix the connecting-map sign by its representative formula.
\item \textbf{Instance and corollaries.} Add the Laurent coefficient factor and the supplied bigrading infrastructure, fill the 22 fields in the table, and instantiate the two existing arrangement equivalences. Implement the worked regressions along the way.
\end{enumerate}

\paragraph{Recommended interfaces inside the new namespace.}
Separate operations on form degrees in $\N$ from zero-extension to $\Z$. The forms themselves are $A$-modules, but $d$, the complexes, their cohomology, and chain homotopies are $k$-linear. Contraction is $A$-linear; Lie derivative is only $k$-linear. The module quotient $\LL/f\LL$ has a differential because its stability was proved, not because $d$ was $A$-linear. These distinctions should be encoded in types.

Use \planned{ExteriorPower} or the homogeneous components of \planned{ExteriorAlgebra} for $\Om_A^n$. The polynomial construction should expose a transport theorem identifying the quotient presentation with these types, so downstream lemmas can use the intrinsic wedge and universal $da$ notation. For morphisms $A\to A'$, the maps on forms are semilinear over the algebra map and ordinary $k$-linear chain maps. They are not generally $A'$-linear maps with the original source module structure.

For the pole proof, keep the explicit identities on $V=\LL/f\LL$ as separate lemmas before invoking any complex-homotopy API. A developer can then package the identity as the desired library homotopy notion, even if its shift convention differs. The finite-reduction lemma avoids any dependency on a filtered-complex spectral sequence. No construction needs a coordinate cover, a chosen dimension of $A$, or a global free basis of $\Om_A^1$.

\section{Mathlib entry points and verification status}

The requested target is the team's Lean 4.36.0-rc2 / Mathlib checkout. No local Lean executable was available here. Every \planned{OrlikSolomon.DeRham} name is a proposal, not an assertion that a declaration exists. The following existing names were verified in the online Mathlib documentation on 6 November 2026; their exact types must still be checked in the team's pinned checkout.

{\small
\begin{longtable}{@{}p{0.57\textwidth}p{0.38\textwidth}@{}}
\toprule
Verified existing declaration or module & Relevance\\\midrule\endhead
\nolinkurl{KaehlerDifferential.D} &
Universal derivation.\\
\nolinkurl{Derivation.liftKaehlerDifferential} &
Universal map from one-forms to a module.\\
\nolinkurl{KaehlerDifferential.mvPolynomialEquiv} and
\nolinkurl{KaehlerDifferential.mvPolynomialBasis} &
Free polynomial differentials used before descent.\\
\nolinkurl{Algebra.FormallySmooth.liftOfSurjective} and
\nolinkurl{Algebra.FormallySmooth.comp_liftOfSurjective} &
Square-zero section and derivation-lifting arguments.\\
\nolinkurl{retractionKerCotangentToTensorEquivSection} &
Existing conormal/section correspondence; optional shortcut for packaging the explicit proof.\\
\nolinkurl{Mathlib.RingTheory.Smooth.Basic} &
Formal smoothness, smoothness, and projectivity of the K\"ahler module.\\
\nolinkurl{Mathlib.RingTheory.Kaehler.Polynomial} &
Polynomial differential descriptions.\\
\nolinkurl{Mathlib.Algebra.Polynomial.Laurent} &
Laurent polynomials and localization away from the polynomial variable.\\
\bottomrule
\end{longtable}
}
Mathlib's formal-smoothness class is expressed through projectivity and cotangent conditions; its documented lifting results supply the square-zero lifting characterization used in the proofs. Thus the proof of projectivity above may be replaced by the existing instance, and the section construction may directly use \planned{liftOfSurjective}. These are existing foundations, not additional cohomology axioms.

The prompt also identifies the exterior-power, Jacobi--Zariski, tensor-product, and homological-algebra libraries. Exact convenience-lemma names in those libraries have not been verified here and are not claimed. In particular a ready-made de Rham differential, logarithmic complex, or polynomial Poincar\'e lemma has not been assumed. Their construction is the content of this blueprint.

\paragraph{Acceptance criteria.}
Check every theorem-like environment has exactly one label, all new labels use the prescribed \planned{dr-} prefix, and all internal dependencies resolve without cycles after treating the listed first-blueprint labels as external. Then implement the declarations, compile against the fixed interface without adding hypotheses to its fields, and inspect the final instance and corollaries for unproved axioms. The theory structure must not contain pole comparison, residue exactness, or affine-line invariance as extra assumptions on a proposed de Rham object: the proofs above supply them. Verify that the unit and zero-ring regressions typecheck, and that the genuinely non-flat example does not require a flatness instance. A successful LaTeX build or dependency-graph check is not Lean verification.

\paragraph{Sources for context and API verification.}
The proofs needed for the construction are included above. The Stacks Project discussion of logarithmic poles provides context for the residue sequence and the smooth-pair hypotheses:
\url{https://stacks.math.columbia.edu/tag/0FMU}.
The Mathlib declarations were checked at:
\begin{itemize}
\item \url{https://leanprover-community.github.io/mathlib4_docs/Mathlib/RingTheory/Kaehler/Basic.html}
\item \url{https://leanprover-community.github.io/mathlib4_docs/Mathlib/RingTheory/Kaehler/Polynomial.html}
\item \url{https://leanprover-community.github.io/mathlib4_docs/Mathlib/RingTheory/Smooth/Basic.html}
\item \url{https://leanprover-community.github.io/mathlib4_docs/Mathlib/RingTheory/Smooth/Kaehler.html}
\item \url{https://leanprover-community.github.io/mathlib4_docs/Mathlib/Algebra/Polynomial/Laurent.html}
\end{itemize}
These references are not substitutes for a proof of the new differential, residue, or pole-comparison constructions.

\end{document}
